% !TeX program = xelatex
% Самодостатній файл UTF-8. Компіляція: xelatex <ім’я-файлу>.tex (двічі).
\documentclass[10pt,a4paper,landscape]{article}
\usepackage[margin=15mm]{geometry}
\usepackage{fontspec}
\setmainfont{FreeSerif.otf}[BoldFont=FreeSerifBold.otf,ItalicFont=FreeSerifItalic.otf,BoldItalicFont=FreeSerifBoldItalic.otf]
\usepackage{polyglossia}
\setdefaultlanguage{ukrainian}
\usepackage{longtable,array,booktabs,xcolor,fancyhdr}
\usepackage[hidelinks]{hyperref}
\definecolor{accent}{HTML}{16656B}
\pagestyle{fancy}\fancyhf{}
\fancyhead[L]{STEM • 5 клас}\fancyhead[R]{Навчальний та календарно-тематичний план}
\fancyfoot[C]{\thepage}
\setlength{\headheight}{14pt}
\setlength{\parindent}{0pt}
\setlength{\parskip}{5pt}
\renewcommand{\arraystretch}{1.18}
\newcolumntype{P}[1]{>{\raggedright\arraybackslash}p{#1}}
\begin{document}
{\LARGE\color{accent} Навчальний план STEM, 5 клас}

{\large 1,5 год/тиждень • 48 годин на рік}

Заклад освіти: \rule{90mm}{0.3pt}\quad Навчальний рік: \rule{30mm}{0.3pt}

Учитель/учителька: \rule{95mm}{0.3pt}\quad Клас: \rule{25mm}{0.3pt}
\section*{Пояснювальна записка}
Авторський адаптований план укладено за наданою модельною програмою «STEM. 5--6 класи (міжгалузевий інтегрований курс)» О. В. Бутурліної та О. Є. Артєм’євої (розділ 5 класу, с. 9--20 PDF) і зошитом-конспектом «STEM-LAB. 5» (О. Бутурліна та ін., за загальною редакцією О. В. Коршунової, ВД «Освіта», 2023; файл «Stem. Full.pdf»).
Мета: формувати дослідницькі, інженерні, математичні й цифрові вміння, навички співпраці та уявлення про STEM-професії через практичні проєкти.

У наданій програмі орієнтовний розподіл становить 35 годин: 3 + 7 + 5 + 5 + 5 + 5 + 5. Програма дозволяє вчителеві змінювати розподіл годин зі збереженням очікуваних результатів (с. 8 PDF). Цей план адаптовано до 32 навчальних тижнів. Збережено вступ, п’ять змістових модулів і підсумкову STEM-практику. Послідовність модулів відповідає програмі; у зошиті розділ про сили передує розділу про Всесвіт.
\textbf{Облік часу:} І семестр — $15\times1{,}5=22{,}5$ год; ІІ семестр — $17\times1{,}5=25{,}5$ год; разом 48 год. Один рядок календарної таблиці — тижневий тематичний блок: 1 година основної роботи та 0,5 години практичного поглиблення. Це навчальні, а не астрономічні години; додатковий час є роботою з учителем, а не домашнім завданням.

Для розкладу з цілими уроками можна чергувати 1 і 2 уроки, отримуючи 3 години за два тижні. Однак за непарних 15 і 17 тижнів таке чергування дає 22 і 26 (або 23 і 25) цілих уроків, а не точні 22,5 і 25,5. Для точного семестрового обліку потрібне погоджене школою блокове планування з обліком півгодинних частин; дати й фактичні заняття вчитель вписує відповідно до розкладу. Півгодинні частини можна об’єднувати в практикуми в межах семестру.

\textbf{Ресурси та безпека:} зошит, папір, картон, лінійки, фарби, ліхтарик, низьковольтні навчальні набори; за можливості — комп’ютер зі Scratch і камера. Цифрові завдання мають паперовий варіант. Електричні моделі працюють лише від батарейок; не використовувати мережеву напругу. Не спрямовувати оптичні прилади на Сонце або очі. Запуски моделей проводити у вільній зоні. Складні вирізи готує вчитель.

\textbf{Джерельні позначення:} «Зошит» у таблиці — номер сторінки наданого 64-сторінкового скану, з обкладинкою як сторінкою 1. Це опора для добору вправ, а не вимога виконати всі завдання сторінок. Частину практикумів, Scratch, семестрову діагностику й хакатон адаптовано або запропоновано за програмою; вони не подаються як дослівні завдання зошита.
\newpage
\section*{Розподіл навчального часу}
\begin{tabular}{P{125mm}rrr}
\toprule Розділ & І семестр & ІІ семестр & Разом \\\midrule
Вступ & 4,5 & 0 & 4,5 \\
Людина — людина. Я у школі & 9 & 0 & 9 \\
Людина — природа. Я у Всесвіті & 9 & 0 & 9 \\
Людина — техніка. Сила — це сила & 0 & 7,5 & 7,5 \\
Людина — образ. Промінь і світло & 0 & 6 & 6 \\
Людина — знак. Під знаком STEM & 0 & 6 & 6 \\
Хакатон, фестиваль, STEM-практика & 0 & 6 & 6 \\
\midrule Разом & 22,5 & 25,5 & 48 \\
\bottomrule\end{tabular}
\section*{Оцінювання та очікувані результати}
На першому занятті — стартова бесіда й короткі практичні завдання без окремої додаткової години. Упродовж навчання — спостереження за роботою, аналіз таблиць, ескізів, моделей, короткі пояснення учнів. Після кожного проєкту — самооцінювання та взаємний відгук. Семестрові підсумки інтегровані у тижні 15 і 32.

Спільні критерії: зрозуміла мета; коректні спостереження й розрахунки; працездатність та безпечність моделі; пояснення змін після випробування; співпраця й особистий внесок. Для формувального відгуку використовуються описи «виконує з допомогою», «виконує самостійно», «пояснює та переносить у нову ситуацію». Конкретну шкалу підсумкового оцінювання визначає заклад освіти.

Наприкінці року учень/учениця формулює просту проблему й гіпотезу, планує послідовність дій, вимірює та подає дані в таблиці або діаграмі, створює і вдосконалює модель, користується знаками та цифровими засобами, пояснює рішення і називає пов’язані професії. Портфоліо містить особистий план, результати п’яти тематичних проєктів і підсумковий прототип.
\newpage
\section*{Календарно-тематичний план. І семестр}
\small\begin{longtable}{P{13mm}P{16mm}P{47mm}P{18mm}P{88mm}P{52mm}}
\toprule Тиждень / год & Дата & Модуль і тема & Зошит & Навчальна діяльність & Очікуваний результат \\\midrule\endfirsthead
\toprule Тиждень / год & Дата & Модуль і тема & Зошит & Навчальна діяльність & Очікуваний результат \\\midrule\endhead
\bottomrule\endfoot
1 / 1,5 & & \textbf{Вступ}\par STEM-освіта та світ професій & 2--5 & \textbf{1 год:} Коло асоціацій STEM; сортування професій; стартова бесіда. \par\textbf{0,5 год:} Скласти діаграму професій за результатами опитування класу. & Пояснює складники STEM і наводить приклади професій. \\ \midrule
2 / 1,5 & & \textbf{Вступ}\par Проєкт: від задуму до плану & 6--9 & \textbf{1 год:} Сформулювати мету шкільного STEM-центру, розподілити ролі, скласти резюме проєкту. \par\textbf{0,5 год:} Порівняти два задуми STEM-центру за доступними ресурсами. & Визначає мету, продукт, ресурси й етапи проєкту. \\ \midrule
3 / 1,5 & & \textbf{Вступ}\par Мій шлях до успіху & 6--7 & \textbf{1 год:} Створити особистий план розвитку; обговорити рівні можливості у виборі професії. \par\textbf{0,5 год:} Додати до портфоліо карту сильних сторін і взаємний відгук. & Визначає власну навчальну ціль і спосіб перевірки поступу. \\ \midrule
4 / 1,5 & & \textbf{Людина — людина. Я у школі}\par Моя школа: люди та професії & 10--12 & \textbf{1 год:} Мініопитування про шкільні професії; таблиця й діаграма складу класу. \par\textbf{0,5 год:} Підготувати запитання для інтерв’ю зі шкільним працівником. & Збирає знеособлені дані й пояснює діаграму. \\ \midrule
5 / 1,5 & & \textbf{Людина — людина. Я у школі}\par Територія школи. Масштаб і план & 13--14 & \textbf{1 год:} Виміряти клас; обчислити площу; виконати план у вибраному масштабі. \par\textbf{0,5 год:} Порівняти план із картою; знайти та виправити похибки вимірювання. & Переносить виміри на план і застосовує умовні позначення. \\ \midrule
6 / 1,5 & & \textbf{Людина — людина. Я у школі}\par Умови та засоби навчання & 15--16 & \textbf{1 год:} Порівняти навчальні засоби різних часів; спланувати дослідження класу. \par\textbf{0,5 год:} Провести командне дослідження освітленості на якісному рівні. & Обґрунтовує вибір засобів навчання за заданими критеріями. \\ \midrule
7 / 1,5 & & \textbf{Людина — людина. Я у школі}\par Чому енергоефективність важлива & 17--19 & \textbf{1 год:} Порівняти споживання ламп за наданими даними; запропонувати способи економії. \par\textbf{0,5 год:} Перевірити розрахунки для іншої тривалості освітлення; побудувати діаграму. & Обчислює витрати енергії за зразком і пояснює вибір рішення. \\ \midrule
8 / 1,5 & & \textbf{Людина — людина. Я у школі}\par Енергоефективна сигнальна система & 20--21 & \textbf{1 год:} Скласти низьковольтне коло з батарейки, вимикача та світлодіодного модуля. \par\textbf{0,5 год:} Випробувати модель; знайти причину несправності та поліпшити з’єднання. & Читає просту схему, збирає і перевіряє модель за інструкцією. \\ \midrule
9 / 1,5 & & \textbf{Людина — людина. Я у школі}\par Школа моєї мрії: захист проєкту & 22--23 & \textbf{1 год:} Зібрати простий макет із підготовлених деталей; представити енергоощадне рішення. \par\textbf{0,5 год:} Доопрацювати макет за відгуками; оформити паспорт проєкту. & Пов’язує план, модель і потреби користувачів; оцінює внесок команди. \\ \midrule
10 / 1,5 & & \textbf{Людина — природа. Я у Всесвіті}\par Об’єкти Всесвіту. Сонячна система & 36--38 & \textbf{1 год:} Побудувати порівняльну модель планет; обговорити масштаб і межі моделі. \par\textbf{0,5 год:} Створити просту анімацію руху об’єкта в Scratch; за відсутності техніки — алгоритм на папері. & Розрізняє основні об’єкти й пояснює умовність їх зображення. \\ \midrule
11 / 1,5 & & \textbf{Людина — природа. Я у Всесвіті}\par Зоряне небо та оптичні прилади & 39--40 & \textbf{1 год:} Виготовити зоряний ліхтарик за шаблоном; порівняти телескоп і бінокль. \par\textbf{0,5 год:} Змінити відстань до екрана; записати, як змінюється зображення. & Пояснює призначення оптичних приладів і демонструє модель. \\ \midrule
12 / 1,5 & & \textbf{Людина — природа. Я у Всесвіті}\par Від повітряної кулі до космічної станції & 41--42 & \textbf{1 год:} Створити стрічку освоєння космосу; змоделювати станцію з готових елементів. \par\textbf{0,5 год:} Розробити схему життєзабезпечення станції й розподіл роботи екіпажу. & Називає функції станції та обов’язки екіпажу. \\ \midrule
13 / 1,5 & & \textbf{Людина — природа. Я у Всесвіті}\par Людина у відкритому космосі & 43--44 & \textbf{1 год:} Порівняти захисні матеріали; спроєктувати модель захисту від холоду та ударів. \par\textbf{0,5 год:} Порівняти два зразки захисту в безпечному модельному випробуванні. & Пов’язує небезпеку з властивістю захисного матеріалу. \\ \midrule
14 / 1,5 & & \textbf{Людина — природа. Я у Всесвіті}\par Космічні професії. Україна космічна & 45--46 & \textbf{1 год:} Підготувати картку професії й інформаційний буклет за матеріалами вчителя. \par\textbf{0,5 год:} Перевірити твердження буклета за другим джерелом; додати ілюстративну схему. & Називає компетентності фахівця; зазначає джерела інформації. \\ \midrule
15 / 1,5 & & \textbf{Людина — природа. Я у Всесвіті}\par Я у Всесвіті: представлення результатів & 36--46 & \textbf{1 год:} Представити космічні моделі та буклети; виконати коротку семестрову діагностику. \par\textbf{0,5 год:} Виправити помилки діагностики; удосконалити модель за відгуками. & Пояснює рішення з опорою на дані; визначає, що варто повторити. \\ \midrule
\end{longtable}\normalsize
\newpage
\section*{Календарно-тематичний план. ІІ семестр}
\small\begin{longtable}{P{13mm}P{16mm}P{47mm}P{18mm}P{88mm}P{52mm}}
\toprule Тиждень / год & Дата & Модуль і тема & Зошит & Навчальна діяльність & Очікуваний результат \\\midrule\endfirsthead
\toprule Тиждень / год & Дата & Модуль і тема & Зошит & Навчальна діяльність & Очікуваний результат \\\midrule\endhead
\bottomrule\endfoot
16 / 1,5 & & \textbf{Людина — техніка. Сила — це сила}\par Сили у природі. Історія повітроплавання & 24--25 & \textbf{1 год:} Дослідити дію сили на рух і форму предмета; порівняти літальні апарати. \par\textbf{0,5 год:} Скласти схему сил, що діють на модель літального апарата. & Наводить приклади дії сил і розрізняє спостереження та припущення. \\ \midrule
17 / 1,5 & & \textbf{Людина — техніка. Сила — це сила}\par Сила тертя & 26--28 & \textbf{1 год:} Порівняти рух однакового тіла на різних поверхнях; заповнити таблицю. \par\textbf{0,5 год:} Повторити вимірювання; пояснити розбіжності та корисне застосування тертя. & Робить висновок про тертя, змінюючи одну умову досліду. \\ \midrule
18 / 1,5 & & \textbf{Людина — техніка. Сила — це сила}\par Прості машини та механізми & 29--31 & \textbf{1 год:} Дослідити важіль на лінійці; знайти прості механізми у побутових предметах. \par\textbf{0,5 год:} Порівняти положення опори й вантажу; оформити схему спостережень. & Пояснює дію важеля на основі досліду. \\ \midrule
19 / 1,5 & & \textbf{Людина — техніка. Сила — це сила}\par Крила: конструювання моделі & 32--34 & \textbf{1 год:} Виготовити паперовий літак; висунути припущення про вплив форми крил. \par\textbf{0,5 год:} Створити другий варіант крил, зберігаючи інші умови. & Створює модель за ескізом і формулює перевірювану гіпотезу. \\ \midrule
20 / 1,5 & & \textbf{Людина — техніка. Сила — це сила}\par Крила: випробування й захист & 32--35 & \textbf{1 год:} Виміряти дальність повторних запусків; порівняти результати та представити висновок. \par\textbf{0,5 год:} Побудувати діаграму запусків; повторно перевірити вдосконалену модель. & Фіксує результати й аргументує поліпшення моделі. \\ \midrule
21 / 1,5 & & \textbf{Людина — образ. Промінь і світло}\par Світло і колір & 47--49 & \textbf{1 год:} Дослідити спектр на готовій моделі спектроскопа; порівняти кольорові зображення. \par\textbf{0,5 год:} Зібрати простий спектроскоп за підготовленим шаблоном і порівняти безпечні джерела світла. & Описує спектральний склад білого світла й результати спостереження. \\ \midrule
22 / 1,5 & & \textbf{Людина — образ. Промінь і світло}\par Фарби та відтінки & 50--51 & \textbf{1 год:} Змішати фарби; створити палітру; порівняти природні барвники. \par\textbf{0,5 год:} Дослідити вплив співвідношення фарб; записати рецепт відтінку. & Отримує відтінки й відрізняє змішування фарб від змішування світла. \\ \midrule
23 / 1,5 & & \textbf{Людина — образ. Промінь і світло}\par Поезія та фотографія: образ Сонця & 52--55 & \textbf{1 год:} Зіставити художній образ із явищем; створити фотографію або паперову композицію світла і тіні. \par\textbf{0,5 год:} Порівняти два варіанти освітлення одного предмета; оформити фотопару. & Добирає засоби зображення й пояснює роль освітлення. \\ \midrule
24 / 1,5 & & \textbf{Людина — образ. Промінь і світло}\par Проєкт «Намалюю тобі Сонце» & 47--55 & \textbf{1 год:} Оформити творчу роботу з палітрою, світлиною та поясненням світлового явища. \par\textbf{0,5 год:} Провести мініекспозицію та вдосконалити роботу за відгуками. & Поєднує наукове пояснення і художній задум; презентує результат. \\ \midrule
25 / 1,5 & & \textbf{Людина — знак. Під знаком STEM}\par Знаки й професії системи «людина — знак» & 56--58 & \textbf{1 год:} Класифікувати знаки; закодувати коротке повідомлення; обговорити професії. \par\textbf{0,5 год:} Обмінятися шифрами; перевірити однозначність декодування. & Розрізняє природні й штучні знаки та пояснює правила кодування. \\ \midrule
26 / 1,5 & & \textbf{Людина — знак. Під знаком STEM}\par Графічні знаки. Геральдика & 57--60 & \textbf{1 год:} Розглянути герби; створити ескіз сімейного герба з поясненням символів. \par\textbf{0,5 год:} Виготовити просту печатку з безпечних готових матеріалів; перевірити відбиток. & Пояснює зв’язок символу зі змістом і читає графічні повідомлення. \\ \midrule
27 / 1,5 & & \textbf{Людина — знак. Під знаком STEM}\par Кодування імені. STEM-символіка & 58--61 & \textbf{1 год:} Закодувати ім’я за обраною абеткою; створити ескіз емблеми. \par\textbf{0,5 год:} Перевірити код у парі; створити другий варіант емблеми. & Користується узгодженим кодом і обґрунтовує символи. \\ \midrule
28 / 1,5 & & \textbf{Людина — знак. Під знаком STEM}\par Проєкт «Я обираю STEM» & 61--63 & \textbf{1 год:} Оформити емблему та презентувати її значення; провести взаємооцінювання. \par\textbf{0,5 год:} Перевірити читабельність емблеми різного розміру; завершити портфоліо. & Передає ідеї STEM символами й пояснює власне рішення. \\ \midrule
29 / 1,5 & & \textbf{Хакатон, фестиваль, STEM-практика}\par STEM-практика: професії громади & 64 & \textbf{1 год:} Провести очну або віртуальну зустріч із фахівцем; за відсутності — дослідити підготовлений кейс. \par\textbf{0,5 год:} Підготувати інформаційну картку місцевої STEM-організації. & Ставить змістовні запитання й пов’язує професію з навчальними вміннями. \\ \midrule
30 / 1,5 & & \textbf{Хакатон, фестиваль, STEM-практика}\par Хакатон: задум і прототип & 62--64 & \textbf{1 год:} Обрати проблему класу; спланувати й створити простий прототип із доступних матеріалів. \par\textbf{0,5 год:} Провести першу перевірку прототипу за критеріями. & Визначає проблему, критерії успіху, ресурси й ролі. \\ \midrule
31 / 1,5 & & \textbf{Хакатон, фестиваль, STEM-практика}\par Хакатон: випробування і вдосконалення & 62--64 & \textbf{1 год:} Випробувати прототип; зафіксувати зміни; підготувати короткий виступ. \par\textbf{0,5 год:} Повторити випробування та порівняти початковий і кінцевий результати. & Обґрунтовує зміни результатами випробування. \\ \midrule
32 / 1,5 & & \textbf{Хакатон, фестиваль, STEM-практика}\par STEM-фестиваль. Підсумкова рефлексія & 62--64 & \textbf{1 год:} Представити портфоліо й прототип; виконати підсумкові завдання та самооцінювання. \par\textbf{0,5 год:} Обговорити індивідуальні результати й сформулювати цілі на наступний рік. & Демонструє поступ у дослідженні, конструюванні та співпраці. \\ \midrule
\end{longtable}\normalsize
\end{document}
